\subsection{Multiplicities over a local unit ring}
\label{subsec:local-unit-projectives}

\begin{proof}[Proof of Theorem~\ref{thm:local-unit-projectives}]
First use Theorem~\ref{thm:integral-support-reduction} to replace $F$ by
its image-ideal diagram. Fix a finite-projective support $L$.
Theorem~\ref{thm:action-colimits} identifies its action category with
the ideal category
\begin{equation}\label{eq:local-ideal-category}
 \mathcal A_L(l,k)=H(l,k):=\bigcup_{ql\leq k}I_q\subseteq R.
\end{equation}
The union is directed by finite joins of the scalar indices. Since the
transition maps are inclusions, this is an actual union, and composition
is the ordinary product in $R$. In particular $H(l,l)=R$, because the
index $1$ is available and is greatest.

The set $W$ contains $1$, is upward closed, and is closed under
multiplication: if $I_q=I_p=R$, then
$R=I_qI_p\subseteq I_{qp}$. Hence $\preceq_W$ is a preorder. Moreover
\[
 H(l,k)=R\quad\Longleftrightarrow\quad l\preceq_W k.
\]
Indeed, the union in \eqref{eq:local-ideal-category} contains $1_R$
exactly when one of its constituent ideals contains $1_R$.

Let $P_l=\mathcal A_L(l,-)$ be the covariant representable. Yoneda gives
$\End(P_l)\simeq R^{\mathrm{op}}$, a nonzero local ring, so $P_l$ is
indecomposable. Two representables $P_l,P_k$ are isomorphic exactly
when there are inverse coefficients in $H(l,k)$ and $H(k,l)$.
Inverse coefficients are units of $R$, so both ideals are $R$.
Conversely, if both ideals are $R$, their coefficients $1_R$ give
inverse maps. Thus
\[
 P_l\simeq P_k\quad\Longleftrightarrow\quad l\sim_W k.
\]

Every finite projective amplitude is a retract of a finite direct sum
$P_{l_1}\oplus\cdots\oplus P_{l_s}$. The ordinary Krull--Schmidt
argument applies even though $\mathcal A_L$ may have infinitely many
objects. To see the finiteness explicitly, take the endomorphism ring
$\Gamma$ of this finite sum. The corresponding decomposition of its
regular module has local endomorphism rings. By
\cite[Proposition~4.1, pp.~8--9]{Krause}, $\Gamma$ is semiperfect and
its finite projectives form a Krull--Schmidt category. The additive
Yoneda equivalence for the direct summands of finite sums identifies
the original retract with one of these projectives. Uniqueness and
the description of direct summands in
\cite[Theorem~4.2 and Corollary~4.3, p.~9]{Krause} show that it is a
finite sum of the same representable isomorphism types, with unique
multiplicities. These locators refer to the author's manuscript. No
Artinian or finite-ring hypothesis has been used.

It remains to impose the canonical ideal condition. For any ideal
$\mathcal J$ of an additive category,
\[
 (\mathcal J P_l)(k)=\mathcal J(l,k).
\]
Composition gives one inclusion, and composition with $1_l$ gives the
other. Therefore $\mathcal J P_l=P_l$ exactly when
$1_l\in\mathcal J(l,l)$. For the present inclusion diagram, the
finite-frame formula gives
\[
 \mathcal J_L(l,l)=\bigcup_{q\in C(l,l)}I_q.
\]
This ideal is $R$ exactly when $C(l,l)$ meets $W$. The intrinsic
description $C(l,l)=\mathord\downarrow\{\psi(l):\psi\in D_l\}$ and
upward closure of $W$ give the equivalent admissible-functional
condition. It is invariant under isomorphism of representables because
$\mathcal J_L$ is two-sided, proving independence of the representative.

The ideal acts componentwise on a direct sum. Consequently
$\mathcal J_LG=G$ exactly when all representable types occurring with
positive multiplicity are allowed. Theorem~\ref{thm:canonical-ideal}
now proves necessity and sufficiency of the stated data.

Formula~\eqref{eq:local-unit-fibers} is the resulting sum of
representables. If $x\in I_p$ and $y\in I_q$ with $ql\leq k$, then
$xy\in I_{pq}$ and $(pq)l\leq pk$, so left multiplication has the
asserted target. The transitions, naturality, associativity and unit
law follow from inclusion and ordinary ring multiplication. Support
isomorphisms preserve all scalar inequalities and transport the
allowed classes; conversely a semimodule isomorphism induces a support
isomorphism and, by the uniqueness just proved, the same multiplicities.

Finally lift the finite idempotent for this separated projective along
the coefficient quotient. The preceding theorem gives its unique
original projective isomorphism class. Restriction of the separated
module along the quotient is not the inverse construction. A zero
multiplicity function gives the zero-amplitude module on $L$, which
is the zero semimodule only if $L$ is itself the zero support.
\end{proof}

If $I_0=R$, all support objects are $\sim_W$-equivalent and allowed,
leaving the ordinary free rank over the local ring together with $L$.
If $I_0$ is proper, singleton support admits no nonzero projective
amplitude. If the unit ring is zero, outside the nonzero-local
hypothesis, the unit law forces every coefficient group to be zero,
so only the projective $Q$-support remains.
